\section{An Information-Theoretic Account}
\label{sec:theory}

We formalise why a coherence-maximising retriever can be paradox-prone.
Consider a query $Q$, a candidate set $\mathcal{C} = \{c_1, \ldots,
c_n\}$, and an unobserved ground-truth answer $Y$. Let
$\textrm{coh}(\mathcal{C})$ denote any monotone-increasing function of
the off-diagonal pairwise similarity of the candidate set --- CCS is
one such function. Let $I(\mathcal{C}; Y \mid Q)$ denote the mutual
information between the retrieved set and the answer, conditioned on
the query.

\begin{proposition}[Refinement Paradox]
\label{prop:paradox}
Let $\mathcal{C}_\star \subseteq \mathcal{C}$ be the subset that
maximises $\textrm{coh}(\cdot)$ subject to $|\mathcal{C}_\star| = k'$
($k' \leq |\mathcal{C}|$). If the per-passage relevance scores
$\{r(c_i, Q)\}_i$ used by the refinement step are not perfectly aligned
with each passage's marginal contribution to $I(\mathcal{C}; Y \mid Q)$,
then there exist queries for which
$I(\mathcal{C}_\star; Y \mid Q) < I(\mathcal{C}; Y \mid Q)$. In words:
maximising coherence can strictly reduce the information content of
the retained set.
\end{proposition}

The proof is a one-line application of the data-processing inequality:
the refinement step is a stochastic map $\mathcal{C} \to
\mathcal{C}_\star$ that depends only on per-passage relevance, so any
information uniquely supplied by a passage that is removed for being
incoherent (with the rest) is lost --- regardless of how relevant that
passage was. The inequality is strict whenever such a passage exists.

\begin{theorem}[Recovery Condition]
\label{thm:paradox_free}
A retriever $\mathcal{R}$ is paradox-free if it satisfies:
$\forall \mathcal{C}, \, I(\mathcal{R}(\mathcal{C}); Y \mid Q) \geq
I(\mathcal{C}; Y \mid Q)$. A sufficient condition is that
$\mathcal{R}$ retain every passage whose marginal information
contribution exceeds a fixed threshold, even when its per-passage
similarity is low.
\end{theorem}

\paragraph{Operationalisation.}
The HCPC-v$2$ rank-protection rule (\S\ref{sec:method}) is one
implementation of the recovery condition: by always retaining the
top-$k$ passages by raw query similarity, we preserve their
information contribution even when refinement would otherwise discard
them. The empirical recovery numbers in
Table~\ref{tab:hcpc_main} are consistent with the bound: the gap
$\textrm{faith}_{\textrm{base}} - \textrm{faith}_{\textrm{v1}}$ tracks
the proportion of queries on which the protected passages would
otherwise have been refined away.

The full information-theoretic derivation, including the dataset-level
prediction (HotpotQA's multi-hop questions exhibit the largest
paradox; SQuAD's single-span answers exhibit the smallest), the
empirical bound verification (respected on $96\%$ of $n=150$
queries), and the connection to selection-bias literature, is in the
longform appendix.
